\subsubsection{差商}\label{subsec:elliptic-regularity-difference-quotients}

设 $\Omega\subset\mathbb{R}^n$, $\{\bm{e}_1,\ldots,\bm{e}_n\}$ 为 $\mathbb{R}^n$ 的标准正交基. 对任意函数 $u\in L^p(\Omega)$, 可以形式地将沿 $\bm{e}_i$ 方向的差商定义为
\begin{equation}\label{eq:coordinate-direction-difference-quotient-definition}
D_i^hu(\bm{x}):=\frac{u(\bm{x}+h\bm{e}_i)-u(\bm{x})}{h}.
\end{equation}
当然, 对靠近边界的 $\bm{x}$, $\bm{x}+h\bm{e}_i$ 可能超出 $\Omega$, 所以这个函数未必对所有 $\bm{x}\in\Omega$ 都有良好定义. 不过可以得到以下结果.

\begin{Lemma}
\label{lem:difference-quotient-weak-derivative-bound}
设 $u\in W^{1,p}(\Omega)$, $1\leq p\leq\infty$. 则对任意 $\Omega^{\prime}\subset\subset\Omega$ 和任意 $h<\operatorname{dist}(\Omega^{\prime},\partial\Omega)$, 有 $D_i^hu\in L^p(\Omega^{\prime})$, 且
\begin{equation}\label{eq:difference-quotient-bound-by-weak-derivative}
\|D_i^hu\|_{L^p(\Omega^{\prime})}\leq\|u_{x_i}\|_{L^p(\Omega)}.
\end{equation}

\end{Lemma}

\begin{Proof}
设 $\Omega^{\prime}$ 和 $h$ 满足引理的假设. 对任意 $\xi\in[0,h]$, 定义
\begin{equation}\label{eq:difference-quotient-translated-interior-domain}
\Omega_{\xi,i}^{\prime}:=\{\bm{x}\in\Omega\mid\bm{x}=\overline{\bm{x}}+\xi\bm{e}_i;\quad\overline{\bm{x}}\in\Omega^{\prime}\}.
\end{equation}
对 $p\in[1,\infty)$, 首先考虑 $u\in C_b^1(\Omega)\cap W^{1,p}(\Omega)$ 的情形. 则对任意 $\bm{x}\in\Omega^{\prime}$, 可利用微积分基本定理写出
\begin{equation}\label{eq:difference-quotient-average-derivative-identity}
\begin{aligned}
D_i^hu(\bm{x})&=\frac{u(\bm{x}+h\bm{e}_i)-u(\bm{x})}{h}\\
&=\frac{1}{h}\int_0^hu_{x_i}(\bm{x}+\xi\bm{e}_i)\,\mathrm{d}\xi.
\end{aligned}
\end{equation}
于是利用 Hölder 不等式并交换积分次序, 得到\[\begin{aligned}
\int_{\Omega^{\prime}}|D_i^hu(\bm{x})|^p\,\mathrm{d}\bm{x}
&\leq\int_{\Omega^{\prime}}\frac{1}{h}\int_0^h|u_{x_i}(\bm{x}+\xi\bm{e}_i)|^p\,\mathrm{d}\xi\,\mathrm{d}\bm{x}\\
&\leq\frac{1}{h}\int_0^h\int_{\Omega_{\xi,i}^{\prime}}|u_{x_i}(\bm{x})|^p\,\mathrm{d}\bm{x}\,\mathrm{d}\xi\\
&\leq\int_\Omega|u_{x_i}|^p\,\mathrm{d}\bm{x}.
\end{aligned}\]由 Theorem~\ref{thm:sobolev-interior-smooth-density}, 取极限便将这个不等式推广到整个空间.

对 $p=\infty$, 注意 \eqref{eq:difference-quotient-average-derivative-identity} 在分布意义下成立. 由于试验函数在 $L^1$ 中稠密, 将 $D_i^hu$ 和 $u_{x_i}$ 视为试验函数集上的线性泛函, 对其进行运算和估计便得该不等式. 这留给读者完成.
\end{Proof}
